\section{Implications for Exact Certification}\label{sec:engineering-corollaries}

The regime hierarchy matters for more than class labels. The same exact question, which coordinates can be omitted without changing the decision boundary, reappears when one tries to simplify models, build reliable certifiers, check witnesses, approximate exact minimizers, or compress a central interface. This section collects those consequences in one place. The common lesson is that exact relevance is expensive to certify and easy to displace rather than remove. One downstream consequence is a trilemma for exact certifiers in the hard regime: once exact certification is hard, every interface-compression story must be read under the possibility of abstention, weakened guarantees, or overclaiming.

\subsection{Exact Simplification and Over-Specification}

\leanmetapending{\LH{CR1}, \LH{CT12}}
\begin{proposition}[Configuration Simplification as \textsc{Sufficiency-Check}]
\label{prop:config-sufficiency}
Let $P=\{p_1,\ldots,p_n\}$ be configuration parameters and let $B$ be a finite set of observable behaviors. Suppose each configuration state determines the subset of behaviors that remain feasible or optimal. Then the question
\[
\text{``Does parameter subset } I \subseteq P \text{ preserve all decision-relevant behavior?''}
\]
is an instance of \textnormal{\textsc{Sufficiency-Check}}.
\end{proposition}


\begin{proof}
Construct a decision problem $\mathcal{D}=(A,X_1,\ldots,X_n,U)$ by taking actions $A=B$, coordinate domains $X_i$ equal to the domains of the parameters $p_i$, and defining $U(b,s)=1$ exactly when behavior $b$ is realized or admissible under configuration state $s$. Then $\Opt(s)$ is the behavior set induced by $s$, and coordinate subset $I$ is sufficient exactly when agreement on the parameters in $I$ forces agreement on the induced optimal-behavior set.
\end{proof}

Worked applied reductions (configuration case study, one-step POMDP reduction, and hyperparameter-redundancy reduction) are moved to Appendix~\ref{app:applications}. This keeps the main theorem spine focused on complexity classification while retaining concrete model translations.


\leanmetapending{\LHrng{DQ}{36}{37}}
\begin{corollary}[No General-Purpose Exact Configuration Minimizer]
\label{cor:no-general-minimizer}
Assuming $\Pclass \neq \coNP$, there is no polynomial-time general-purpose procedure that takes an arbitrary succinctly encoded configuration problem and always returns a smallest behavior-preserving parameter subset.
\end{corollary}


\begin{proof}
Such a procedure would solve \textsc{Minimum-Sufficient-Set} in polynomial time for arbitrary succinctly encoded instances, contradicting Theorem~\ref{thm:minsuff-conp} under the assumption $\Pclass \neq \coNP$.
\end{proof}

\leanmetapending{\LH{DQ26}, \LH{HD14}}
\begin{proposition}[Model-Conditional Rationality of Over-Specification]
\label{prop:rational-overspecification}
Consider a design process in which:
\begin{enumerate}
\item removing a parameter requires certifying that it is irrelevant,
\item exact irrelevance certification is charged according to the computational model of Section~\ref{sec:encoding}, and
\item carrying an extra parameter incurs only linear or otherwise low-order maintenance overhead.
\end{enumerate}
Then, for sufficiently expressive succinctly encoded instances, retaining extra parameters can minimize total cost under the stated cost model.
\end{proposition}


\begin{proof}
By Theorems~\ref{thm:sufficiency-conp} and~\ref{thm:minsuff-conp}, exact sufficiency certification and exact minimization are coNP-complete in the general succinct regime. By Section~\ref{sec:dichotomy}, this hardness is accompanied by exponential lower bounds under ETH in the linear-support regime. If the cost of carrying extra parameters grows only linearly while exact minimization inherits worst-case exponential cost, then beyond a problem-dependent threshold the latter dominates the former. Under that cost model, retaining extra parameters minimizes total cost.
\end{proof}

\begin{remark}
This does not rule out useful simplification tools. It rules out a fully general exact minimizer with worst-case polynomial guarantees. The structured regimes isolated in Section~\ref{sec:tractable} remain viable precisely because they restrict the source of hardness.
\end{remark}

\subsection{Regime-Limited Exact Certification}

\leanmetapending{\LH{DQ36}, \LH{DQ40}, \LH{DQ44}}
\begin{proposition}[Exact Certification Competence Depends on Regime]
\label{prop:competence-by-regime}
Within the model of this paper, exact certification competence is regime-dependent. In the general succinct regime, exact relevance certification and exact minimization inherit the hardness results of Sections~\ref{sec:hardness} and~\ref{sec:dichotomy}. In the structured regimes of Section~\ref{sec:tractable}, exact certification becomes available in polynomial time.
\end{proposition}


\begin{proof}
The negative side is given by Theorems~\ref{thm:sufficiency-conp} and~\ref{thm:minsuff-conp}, together with the ETH-conditioned lower bounds summarized in Section~\ref{sec:dichotomy}. The positive side is given by the tractability results of Section~\ref{sec:tractable}, which provide explicit polynomial-time certification procedures under structural restrictions.
\end{proof}

The next three subsections sharpen that basic regime distinction in different directions: exact reliability under polynomial budgets, witness-checking and approximation limits, and the cost of compressing a central interface while leaving exact relevance unresolved.
